\chapter{Overall Architecture Design}
\label{ch:architecture}

% DRAFT BRIEF: system-level mental model before Ch.5 module detail. Stay at block-diagram
% resolution. Format = block diagrams + 1-2 paragraphs per block.
% Sources: specs 11/10, improvements.md, CLAUDE.md block diagram, i3c_pkg.sv, controller_pkg.sv.

\section{Three-layer architecture}
\brief{Present the design as three layers: (1) PHY (\texttt{i3c_phy} --- 2FF metastability
sync + OD/PP output drivers); (2) controller/protocol (\texttt{controller_active} wrapping
\texttt{flow_active}, the bus serializers, \texttt{scl_generator}, \texttt{bus_monitor},
ENTDAA); (3) host/register (\texttt{csr_register} + four sync FIFOs in \texttt{hci_queues}).
Explain separation of concerns and the design philosophy (simplify, one clock domain, FSM as
single command authority).}

\figph{Three-layer architecture}
  {PHY, controller_active, CSR+HCI layers and their boundaries}{conceptual}{TikZ}

\section{Top module and block diagram}
\brief{\texttt{i3c_controller_top} instantiates \texttt{u_csr} (csr_register), \texttt{u_queues}
(hci_queues: CMD64 + TX/RX/RESP 32), \texttt{u_ctrl} (controller_active and children),
\texttt{u_phy} (i3c_phy). Show the full instance hierarchy and top-level port groups.}

\figph{Top-level block diagram}
  {i3c_controller_top to \{csr, hci_queues, controller_active (flow_active, scl_gen,
   bus_tx_flow/bus_tx, bus_rx_flow, bus_monitor, entdaa_controller/entdaa_fsm), phy\}}
  {src/rtl/i3c_controller_top.sv}{TikZ}

\section{Transaction dataflow}
\brief{Narrate the canonical write and read paths. Write: SW to CSR CMD staging (DWORD0 then
DWORD1 to 64-bit push) to CMD FIFO to \texttt{flow_active} pop to DAT fetch to
\texttt{bus_tx_flow}+\texttt{scl_generator} via \texttt{controller_active} to \texttt{i3c_phy}
to pads to RESP FIFO. Read: same until \texttt{flow_active} enables \texttt{bus_rx_flow}, which
samples SDA on SCL posedge, packs bytes into 32-bit words to RX FIFO to SW read.}

\figph{Transaction dataflow (write + read)}
  {Host to CMD FIFO to flow_active to bus_tx/rx to PHY to pads to RESP/RX, both directions}
  {conceptual}{TikZ}

\section{Clock/reset and signal conventions}
\brief{Single clock domain; async-assert active-low \texttt{rst_ni}; 2FF sync in
\texttt{i3c_phy} (+1 capture flop in \texttt{bus_monitor} = 3 flops before edge detection);
synchronous SW-reset flushing FIFOs + CMD staging. Restate conventions: \texttt{_i}/\texttt{_o},
\texttt{*_valid_i}/\texttt{*_ready_o}, active-low open-drain \texttt{scl_o}/\texttt{sda_o},
\texttt{sel_od_pp_o} (1 = push-pull data, 0 = open-drain addr/ACK).}

\figph{Clock/reset and 2FF sync}
  {single clock domain, async-low reset, 2FF input synchronizer chain}{conceptual}{TikZ}

\section{CHIPS-Alliance comparison (\texorpdfstring{$\sim$}{~}92\% LoC reduction)}
\brief{The simplification story: $\sim$25k to $\sim$2k lines. Per-subsystem comparison
(hand-written CSR vs PeakRDL auto-gen; 32-entry vs larger DAT; no IBI/HDR/multi-master
machinery; folded restart). Attribute the reference (commit SHA) here and in captions.}

\tabph{Per-subsystem LoC comparison}{improvements.md}

\section{FIFO/DAT/descriptor formats}
\brief{Summarise (full bitfields to Appendix A/B): four sync FIFOs (CMD 64-bit, TX/RX/RESP
32-bit), 32-entry DAT (\texttt{dat_entry_t}), the CMD descriptor variants and the RESP
descriptor. Keep to format-overview tables here.}

\tabph{FIFO / DAT / descriptor format overview}{i3c_pkg.sv, controller_pkg.sv}

\section{Error-handling model}
\brief{How errors flow: detection points to \texttt{flow_active} consolidation to RESP
descriptor \texttt{err_status} field. Reference the generated error-code set
(0x0/0x4/0x5/0x6/0x7/0x8/0xA).}

\tabph{Error-status encoding}{phase1_spec_v2.md}
